\documentclass[11pt]{article}
\usepackage[margin=1in,top=0.9in,bottom=0.9in]{geometry}
\usepackage{amsmath,amssymb,amsthm,hyperref}
\hypersetup{colorlinks=true,linkcolor=blue,urlcolor=blue,citecolor=blue}
\setlength{\emergencystretch}{2em}
\newtheorem{theorem}{Theorem}
\newtheorem{lemma}[theorem]{Lemma}
\title{Nonconvergence of randomized subtraction games\\with odd moves}
\author{Research proof and formalization packet}
\date{October 26, 2026}
\begin{document}
\maketitle
\begin{abstract}
Let a finite nonempty subtraction set consist entirely of positive odd
integers, and choose uniformly among the legal moves at each heap.
We prove that the probability of winning for the player to move does not
converge as the heap grows. This establishes the general all-odd
nonconvergence assertion left open by Capitelli and Somma. An alternating
inequality for consecutive probabilities gives positive initial data for
a conserved finite-window sum. Convergence would force this sum to vanish.
The proof requires neither coprime moves nor simplicity of characteristic
roots. The complete recurrence, actual recursively defined probability
sequence, positive invariant, and nonconvergence theorem are verified in Lean.
\end{abstract}

\section{The game and the source question}
Fix a finite nonempty set $S$ of positive integers. From a heap of $n$
chips, a legal move subtracts an element of
$E_n=\{k\in S:k\leq n\}$. The next move is chosen uniformly from $E_n$;
a player with no legal move loses. Let $a_n$ denote the probability that
the player to move wins. Thus
\begin{equation}\label{eq:rec}
 a_n=\begin{cases}
 0,&E_n=\varnothing,\\
 1-\dfrac{1}{|E_n|}\displaystyle\sum_{k\in E_n}a_{n-k},
       &E_n\ne\varnothing.
 \end{cases}
\end{equation}
Positive moves make this a uniquely determined recursion; in particular
$a_0=0$ and $0\leq a_n\leq1$ by induction.

Capitelli and Somma \cite{CS} analyze this exact recurrence. They settle
sets of cardinality below three and prove convergence in their coprime
setting when an even move is present. At the start of Section~3 they
explicitly leave general all-odd nonconvergence open; Question~3.4 asks
about the nonzero coefficient of the root $-1$ in that setting. We give a
direct proof of the missing nonconvergence assertion.

\begin{theorem}\label{thm:main}
For every finite nonempty set $S$ of positive odd integers, the sequence
$(a_n)$ defined by \eqref{eq:rec} has no real limit.
\end{theorem}

This is one universal source-question resolution. The singleton and
two-move cases are included for a complete statement and receive no
separate novelty claim. The neighboring Conjecture~3.3, indexed as
MathDB~363319, concerns simplicity of all characteristic roots. The
present theorem does not assert or require that conjecture.

\section{Consecutive-probability inequalities}
\begin{lemma}\label{lem:zigzag}
If every move in $S$ is odd, then for every $m\geq0$,
\[
 a_{2m}\leq a_{2m+1},\qquad a_{2m+2}\leq a_{2m+1}.
\]
\end{lemma}
\begin{proof}
Prove the inequalities successively by induction. The first begins with
$a_0=0\leq a_1$. At heaps $2m+1$ and $2m+2$ the legal sets coincide,
because no even move can become newly legal. Each odd subtraction gives
an earlier even heap and its following odd heap. The already proved first
inequalities make the latter follower probability at least the former.
Subtracting equal-sized averages in \eqref{eq:rec} gives
$a_{2m+2}\leq a_{2m+1}$. If the legal set is empty, both probabilities
are zero.

Next compare heaps $2m+2$ and $2m+3$. For every old legal move the two
followers are an earlier odd heap and its following even heap. Their
probabilities satisfy the just-established second inequalities, so the
new follower sum is no larger. The only possible newly legal move is
$2m+3$ itself; its follower has probability $a_0=0$. Adding this zero
cannot increase the average. Hence $a_{2m+2}\leq a_{2m+3}$.
If the old legal set is empty, its probability is zero and the new
probability is nonnegative. This covers the initial empty-set branch and
completes the induction.
\end{proof}

Put
\[
 c_n=(-1)^n(1-2a_n),\qquad C(h)=\sum_{j=0}^{h-1}c_j,
 \qquad C(0)=0.
\]
Since every $k\in E_n$ is odd, multiplying the balance form of
\eqref{eq:rec} by $(-1)^n$ gives
\begin{equation}\label{eq:center}
 |E_n|c_n=\sum_{k\in E_n}c_{n-k}\qquad(n\geq0).
\end{equation}
It holds also for empty legal sets as $0=0$. For odd $h=2q+1$,
Lemma~\ref{lem:zigzag} gives
\begin{equation}\label{eq:prefix}
 C(h)=1+2\sum_{j=1}^{q}(a_{2j-1}-a_{2j})\geq1.
\end{equation}

\section{A positive conserved sum}
Let $D=\max S$ and $t=|S|$. Define, for $n\geq0$,
\begin{equation}\label{eq:inv}
 I_n=\sum_{k\in S}\bigl(C(n+D)-C(n+D-k)\bigr)
     =\sum_{k\in S}\sum_{j=0}^{k-1}c_{n+D-k+j}.
\end{equation}

\begin{lemma}\label{lem:invariant}
For every $n\geq0$,
\[
 I_n=I_0=\sum_{k\in S}C(k)\geq t>0.
\]
\end{lemma}
\begin{proof}
Every move is legal at heap $n+D$. Subtracting consecutive windows and
using \eqref{eq:center} gives
\[
 I_{n+1}-I_n=t c_{n+D}-\sum_{k\in S}c_{n+D-k}=0.
\]
It remains to compute the initial value. Sum \eqref{eq:center} over
$0\leq n<D$ and exchange finite sums. The left-hand side becomes
\[
 \sum_{k\in S}\sum_{n=k}^{D-1}c_n
       =\sum_{k\in S}\bigl(C(D)-C(k)\bigr),
\]
whereas the right-hand side is
\[
 \sum_{k\in S}\sum_{n=k}^{D-1}c_{n-k}
       =\sum_{k\in S}C(D-k).
\]
Rearranging gives
\[
 I_0=\sum_{k\in S}\bigl(C(D)-C(D-k)\bigr)
     =\sum_{k\in S}C(k).
\]
The equality is for the total sums; individual suffix and prefix windows
need not agree. Since every $k\in S$ is odd, \eqref{eq:prefix} yields
$I_0\geq |S|>0$.
\end{proof}

\begin{proof}[Proof of Theorem~\ref{thm:main}]
Suppose $a_n\to L$. For $n\geq D$, \eqref{eq:rec} uses the full set $S$.
Taking limits in its finite average gives $L=1-L$, hence $L=1/2$.
Consequently $|c_n|=|1-2a_n|\to0$.
Each of the finitely many shifted terms in the second expression
\eqref{eq:inv} then tends to zero, so $I_n\to0$.
This contradicts Lemma~\ref{lem:invariant}.
\end{proof}

\section{Formal scope, prior work, and novelty limits}
The complete Lean endpoint is
\texttt{OddSubtraction.all\_odd\_game\_nonconvergence}. It quantifies over
an arbitrary finite set, requires nonemptiness and odd positive moves,
and proves that there is no real limit for the actual recursive
\texttt{gameWin} sequence. No supplied probability model, invariant,
spectral result, gcd condition, or root-simplicity hypothesis remains.
The second main endpoint proves the quantitative invariant bound
$I_n\geq|S|$. The packet contains six modules and a pinned isolated replay
of all 37 explicitly inventoried declarations, including the model
structure. The only reported axioms are \texttt{propext},
\texttt{Classical.choice}, and \texttt{Quot.sound}.

The earlier noisy-channel game framework \cite{CP} supplies related
recurrences and fairness questions, but not this all-odd nonconvergence
theorem. Work on ultimately bipartite subtraction games \cite{CH}
concerns optimal-play nim values and related parity phenomena; the
present probabilities average all legal moves rather than select optimal
strategies. The elementary finite-window conservation argument is a tool;
no separate discovery is claimed for it or the known small cases.

Independent review checked the exact primary game definition, all
initial branches, finite-sum indices, nonconvergence argument and final
formal endpoint. Searches of the source title, identifier, author pair
and related randomized-subtraction literature found no earlier full
resolution. The current source record has one version and no journal
reference. These are documented checks, not an absolute certification of
priority; unindexed or unpublished work may remain.

Our formal certificate proves nonconvergence directly. It does not
formalize characteristic-root asymptotics or assert that all roots are
simple. In particular, it must not be described as a proof of the
different MathDB~363319 statement. The source's proposed spectral route
is bypassed by the positive invariant.

\begin{thebibliography}{9}
\bibitem{CS} N. Capitelli and F. Somma,
\emph{On sequences arising from randomizing subtraction games},
arXiv:2405.19593v1 (2024), Sections~1, 3, and 4.1.
\href{https://arxiv.org/abs/2405.19593}{Primary source}.
\bibitem{CP} N. A. Capitelli and M. L. Privitelli,
\emph{Playing through a noisy channel (and knowing it)},
arXiv:2012.06768v1 (2020).
\href{https://arxiv.org/abs/2012.06768}{Primary source}.
\bibitem{CH} G. Cairns and N. B. Ho,
\emph{Ultimately bipartite subtraction games},
Australasian J. Combinatorics 48 (2010), 213--220.
\href{https://ajc.maths.uq.edu.au/pdf/48/ajc_v48_p213.pdf}{Primary source}.
\end{thebibliography}
\end{document}
